\section{Discussion, Limitations, and Open Problems}
\label{sec:discussion}

\subsection{Interpretation}

The survey points to an underdeveloped research object: preserving assumptions
and evidence between existing planners, monitors, and safety mechanisms.  The
edge-contract framework makes those transfers explicit.  Descriptive metadata
becomes contractual when it is bound to assumptions, guarantees, admission
decisions, evidence, and recovery behavior.

The local condition in \cref{eq:contract-composition} does not prove a global
mission property.  Feedback loops, shared variables, circular assumptions,
realizability, probabilistic dependence, lifecycle ordering, and authority
transfer require an implementation-specific formal semantics.  The reference
framework is therefore compatible with, but not a substitute for, the mission,
component, dynamical, and runtime-assurance theories surveyed earlier.

Within the reviewed corpus, the coordinated-autonomy work of Chaki, Dolan,
Giampapa, Porter, and colleagues is the closest mission-level predecessor
because it explicitly relates robot coordination to quantitative assurance and
reliability
\cite{chaki2013assurance,chaki2013probabilistic,porter2013reliability,
stancliff2006missionreliability}.  The edge-contract framework addresses a
different level of abstraction: it inventories the interface-carried
obligations required before a mission-level model, monitor, or reliability
estimate can be applied consistently to a heterogeneous deployed stack.  A
future instantiation should connect the two by making contract rejection,
recovery, and evidence lineage observable variables in a quantitative mission
model.

\emph{Practitioner takeaway.}  A ROS~2 developer can apply the framework by
auditing each inter-robot topic, service, or action.  The interface should
preserve event time separately from reception time, identify frames and units,
define the meaning of uncertainty, and include the relevant configuration,
implementation, and team-epoch identifiers.  The consumer should reject stale
or incompatible inputs and invoke a declared recovery action instead of
continuing silently.  Not every message needs every field; the required
information follows from the downstream component's assumptions
\cite{macenski2022ros2,omg2015dds,kronauer2021latency}.

\subsection{Review Limitations}

The evidence combines formal analyses, simulations, standards, reviews, and
robotics demonstrations created for different purposes.  There is no shared
benchmark from which to estimate prevalence, effect size, or comparative
reliability.  The search was not systematic, and terminology
varies across formal methods, robotics, distributed systems, and embodied AI.
The evidence map can expose coverage and evidence boundaries, but it cannot show
that the proposed framework improves mission outcomes.

The survey also emphasizes the command-to-action path.  It does not attempt a
complete treatment of security, privacy, human factors, certification, or every
multi-robot planning algorithm.  These topics can introduce additional contract
obligations.  Security intersects the proposed identity, provenance, and
authority fields but requires a separate threat model.  Resource availability
is represented as an operational assumption; resource-contract theory is not
surveyed in depth.

\subsection{Open Problems}

Four problems are especially important.  First, the reference fields need
formal semantics that compose across discrete mission logic, continuous
dynamics, probabilistic predictions, and asynchronous communication.  Second,
lifecycle contracts must address membership, partition recovery, ownership, and
retained state without assuming a globally consistent instant.  Third,
uncertainty and evidence need machine-checkable types so that a downstream
consumer can determine whether a proof, calibration result, or deployment test
supports its decision.  Fourth, monitoring must scale without turning every
peer exchange into an unbounded evidence log or centralized bottleneck.

Evaluation should test the framework as an executable mechanism; counting
populated metadata fields is insufficient.  Relevant outcomes include false
admission and rejection, time to detection, correctness of fault attribution, recovery
time, task completion, unsafe events, and runtime overhead.  Reproducible studies
should report the team size, platform, software and middleware versions,
configuration, timing assumptions, injected faults, and run-level traces.  Such
work is necessary before the reference framework can support a reliability or
safety claim.
